\documentclass[11pt,a4paper]{article}
\usepackage[utf8]{inputenc}
\usepackage{amsmath,amssymb,amsfonts}
\usepackage{graphicx}
\usepackage{booktabs}
\usepackage{multirow}
\usepackage{listings}
\usepackage{hyperref}
\usepackage{geometry}
\usepackage{microtype}

\geometry{
    a4paper,
    total={170mm,257mm},
    left=20mm,
    top=20mm,
}

% Code Listings Configuration
\lstset{
    frame=tb,
    language=Python,
    aboveskip=2mm,
    belowskip=2mm,
    showstringspaces=false,
    columns=flexible,
    basicstyle={\small\ttfamily},
    numbers=none,
    breaklines=true,
    breakatwhitespace=true,
    tabsize=4
}

% Hyperref Configuration
\hypersetup{
    colorlinks=true,
    linkcolor=blue,
    filecolor=magenta,
    urlcolor=cyan,
    pdftitle={ShuffleStorm: Distributed Data Management Group Project},
    pdfpagemode=FullScreen,
}

\title{\textbf{ShuffleStorm: Distributed Data Management Group Project}}
\author{\textbf{Georgios Atmatzidis, Vassilis Papazisis} \\ M.Sc.\ Web \& Data Science, Aristotle University of Thessaloniki \\ AEM: 234, 228}
\date{May 2026}

\begin{document}

\maketitle

% ============================================================
%  PART 1 — ALL-PAIRS MATCHING
% ============================================================
\part*{Part 1 — All-Pairs Matching}

\section*{Design and System Architecture}
The system \texttt{shufflestorm} separates dataset ingestion, preprocessing, matching algorithms, and metrics collection into distinct modules. The driver program is \texttt{allpairs.py}; an interactive version is also provided as \texttt{shufflestorm.ipynb}, a Jupyter notebook that reproduces the full pipeline — configuring settings, running all matchers, and displaying collected metrics — without leaving the browser.
\begin{itemize}
    \item \textbf{Preprocessing (\texttt{preprocessing/dataset\_preprocessor.py}):} Extracts relevant columns, projects them to a uniform \texttt{(id, name)} schema, and serialises the processed datasets to CSV.
    \item \textbf{Configuration (\texttt{config/settings.py}):} Centralises dataset names, sizes, group sizes, and the prime $p$ for Afrati-Ullmann.
    \item \textbf{Matchers (\texttt{*\_matcher.py}):} Each strategy is decoupled into its own file with a standardised functional signature.
    \item \textbf{Driver (\texttt{allpairs.py}):} Spins up a local Spark session, loads CSV files into RDDs/DataFrames, runs each matcher, collects metrics (time, task count, shuffle bytes), clears the cache between runs, and saves results to JSON.
\end{itemize}

\section*{Implementation Methods}

\subsection*{Naive Matcher}
The Naive Matcher computes a full Cartesian product via the RDD \texttt{cartesian} transformation, then filters on ID ordering ($id_1 < id_2$) to produce every unique unordered pair exactly once:
\begin{lstlisting}[caption=Naive Matcher]
def run_naive_matching(data_rdd: RDD) -> RDD:
    pairs = data_rdd.cartesian(data_rdd)
    return pairs.filter(lambda x: x[0][0] < x[1][0])
\end{lstlisting}
This generates $\binom{N}{2} = \frac{N(N-1)}{2}$ pairs, making communication cost the dominant bottleneck.

\subsection*{Group-Based Matcher}
The dataset of size $d$ is divided into $g = \lceil d/q \rceil$ groups of capacity $q$. Each record belonging to group $g_i$ is mapped to $g-1$ reducer keys of the form $(\text{sorted}(g_i, g_j), \text{record})$ for every $g_j \neq g_i$, yielding replication rate $r = g-1$. Each reducer then performs an intra-group nested-loop comparison, emitting all pairs — no name-equality filter is applied. A final \texttt{distinct()} removes duplicates produced by overlapping assignments. Replication scales linearly with $g$, causing high shuffle volume for large datasets.

\subsection*{SQL Matcher}
The preprocessed DataFrame is registered as a temporary view and queried with a pure Cartesian product:
\begin{lstlisting}[caption=SQL Matcher Query]
SELECT d1.id, d1.name, d2.id, d2.name
FROM datasets d1 CROSS JOIN datasets d2
WHERE d1.id < d2.id
\end{lstlisting}
There is no equality predicate on a data column, so Spark's Catalyst optimizer cannot rewrite this as an equi-join (e.g.\ \texttt{BroadcastHashJoin} or \texttt{SortMergeJoin}). Instead, Catalyst executes a true Cartesian product — the same $O(N^2)$ plan as the Naive RDD matcher — placing all strategies on equal footing in terms of output semantics. The inequality \texttt{d1.id < d2.id} is evaluated as an inlined post-join filter.

\subsection*{Afrati-Ullmann Matcher}
The dataset is partitioned into $p^2$ cells arranged in a $p \times p$ grid ($p$ prime, $p^2 \le d$). Each record in cell $(i, j)$ is routed to exactly $p+1$ reducers:
\begin{itemize}
    \item For teams $0 \le k < p$: reducer key $(k,\; (i + k \cdot j) \bmod p)$.
    \item For team $p$: reducer key $(p,\; j)$.
\end{itemize}
This guarantees every pair of grid cells meets in at least one reducer. The replication rate $r = p+1$ is theoretically optimal for similarity joins under the MapReduce model. A final \texttt{distinct()} removes duplicates from intra-cell redundant comparisons.

\section*{Evaluation and Comparison}
Benchmarks were run on a single machine in Spark local mode (\texttt{local[*]}). Datasets are from the \href{https://sites.google.com/site/anhaidgroup/useful-stuff/the-magellan-data-repository}{Magellan Data Repository}.

\begin{table}[htbp]
\centering
\caption{Performance Benchmarks — All-Pairs Matching}
\label{tab:results}
\begin{tabular}{llrrrr}
\toprule
\textbf{Dataset} & \textbf{Algorithm} & \textbf{Pairs Found} & \textbf{Time (s)} & \textbf{Tasks} & \textbf{Shuffle (B)} \\
\midrule
\textbf{bikedekho-small}
 & Naive          & 3,405   & 0.742  & 1 & 0          \\
 ($N=500$)
 & Group          & 3,405   & 0.365  & 3 & 63,240     \\
 & SQL            & 3,405   & 0.379  & 4 & 118        \\
 & Afrati-Ullmann & 3,405   & 0.229  & 3 & 72,254     \\
\midrule
\textbf{bikedekho}
 & Naive          & 269,121 & 8.596  & 1 & 0          \\
 ($N=4{,}786$)
 & Group          & 269,121 & 5.750  & 3 & 4,247,456  \\
 & SQL            & 269,121 & 0.458  & 4 & 118        \\
 & Afrati-Ullmann & 269,121 & 4.127  & 3 & 3,435,862  \\
\midrule
\textbf{bikewale}
 & Naive          & 636,262 & 27.135 & 1 & 0          \\
 ($N=9{,}003$)
 & Group          & 636,262 & 18.718 & 3 & 11,089,092 \\
 & SQL            & 636,262 & 0.408  & 4 & 118        \\
 & Afrati-Ullmann & 636,262 & 13.068 & 3 & 7,785,412  \\
\midrule
\textbf{citeseer}
 & SQL    & 2,071,368           & 1.316 & N/A & N/A \\
 ($N=200{,}000$)
 & Others & \textit{Timeout/OOM} & ---   & --- & ---  \\
\bottomrule
\end{tabular}
\end{table}

\noindent\textbf{Naive} scales quadratically (0.74 s at $N=500$ $\to$ 27.14 s at $N=9{,}003$), confirming its $O(N^2)$ cost. \textbf{Afrati-Ullmann} consistently beats Group-Based (13.07 s vs.\ 18.72 s on \texttt{bikewale}) thanks to its grid-based routing which minimises replication to $r = p+1 = 4$; it also produces less shuffle volume (7.8\,MB vs.\ 11.1\,MB on \texttt{bikewale}). \textbf{SQL}, without any equality predicate, forces Catalyst to execute a true Cartesian product and is therefore on the same asymptotic footing as the Naive matcher.

The root cause of the RDD overhead is the \textbf{Python–JVM serialisation boundary}: RDD lambdas force Spark to shuttle data between the JVM and a Python daemon over local sockets on every record. For $O(N^2)$ mappers this overhead dwarfs theoretical improvements. The SQL/DataFrame API stays entirely within the JVM — even when executing a Cartesian product — which is why it outperforms the RDD strategies despite performing the same logical computation: there is no per-record Python deserialisation penalty.

% ============================================================
%  PART 2 — TRIPLE JOIN
% ============================================================
\part*{Part 2: Triple Join: $A(x,y) \bowtie B(y,z) \bowtie C(z,w)$}

The driver is \texttt{triplejoin.py}; as in Part~1, an equivalent Jupyter notebook cell is provided in \texttt{shufflestorm.ipynb} for interactive execution and result inspection.

\textbf{Relations:} Relations are generated synthetically by \texttt{shufflestorm/preprocessing/synthetic\_data \_generator.py} with a configurable size. The synthetic generator produces relations with controlled join selectivity: values for the shared columns $y$ and $z$ are drawn from a small integer domain (size $\approx \sqrt{N}$), so the join output is dense, yielding $\approx 0.099 \cdot N$ output rows per input row on average (confirmed by the row counts in Tables~\ref{tab:triple_relation} and \ref{tab:triple_reducers}).

\textbf{Implementation methods.} Three strategies are implemented in \texttt{shufflestorm/join/} and called from \texttt{triplejoin.py}. All benchmarks are run in \texttt{local[*]} mode with Adaptive Query Execution \emph{disabled} (\texttt{spark.sql.adaptive.enabled = false}) so that the configured shuffle partition count is enforced, keeping results across configurations directly comparable.

\subsection*{Direct Ternary Join (\texttt{ternary\_join.py})}

The direct ternary join implements an algorithm as shown in the lectures for the three-way join, completing $A \bowtie B \bowtie C$ in a \emph{single} distributed pass. We instantiate a $b \times c$ grid of reducers parameterised by two integers $b$ (buckets for $y$) and $c$ (buckets for $z$), giving $b \cdot c$ logical reducers in total.

\medskip\noindent\textbf{Map phase.} Two modular hash functions partition the join-key space:
\[
    h(y) = y \bmod b, \qquad g(z) = z \bmod c.
\]
Each relation is replicated according to how many join attributes it covers:
\begin{itemize}
    \item \textbf{Tuples of $A(x,y)$:} only involve $y$, so they are hashed to bucket $i = h(y)$ and broadcast across all $c$ values of $j$. Each $A$-tuple is replicated $c$ times.
    \item \textbf{Tuples of $B(y,z)$:} span both join attributes and are routed to exactly one reducer $(h(y),\, g(z))$. No replication.
    \item \textbf{Tuples of $C(z,w)$:} only involve $z$, so they are hashed to bucket $j = g(z)$ and broadcast across all $b$ values of $i$. Each $C$-tuple is replicated $b$ times.
\end{itemize}

\medskip\noindent\textbf{Shuffle phase.} The three mapped relations are unioned and grouped by /reducer key $(i,j)$ via \texttt{groupByKey}. A single network pass co-locates all tagged tuples for each reducer.

\medskip\noindent\textbf{Reduce phase.} Each reducer $(i,j)$ separates its tuples by tag, builds hash indices on $B$ (keyed on $y$) and $C$ (keyed on $z$), then executes a local three-way join.

\medskip\noindent\textbf{Replication analysis.} Total shuffle volume scales as $|A| \cdot c + |B| + |C| \cdot b$. With $b = c$ and equal-sized relations of size $N$, this equals $N(2b + 1)$, so the ternary shuffle grows almost linearly with the chosen reducer count, a fact that is visible in Table~\ref{tab:triple_reducers}.

\subsection*{Two Consecutive Binary Joins (\texttt{binary\_join.py})}

The binary strategy decomposes the three-way join into two sequential hash joins at the RDD level, using a reusable helper \texttt{\_hash\_join} function on \texttt{rdd.cogroup}:

\begin{enumerate}
    \item \textbf{Step 1: $A \bowtie_y B$:} Both relations are keyed on $y$. \texttt{cogroup} co-locates all $A$ and $B$ rows sharing the same $y$ onto the same partition. A local nested-loop produces intermediate $AB(x,y,z)$ tuples.
    \item \textbf{Step 2: $AB \bowtie_z C$:} The intermediate $AB$ is keyed on $z$ and joined with $C$ via a second \texttt{cogroup} shuffle. The local nested-loop assembles the final \texttt{Row(x, y, z, w)}.
\end{enumerate}

The key cost is that $AB$ is \emph{fully materialised} before the second shuffle. In our synthetic benchmark the join on $y$ has high fan-out (domain size $\approx \sqrt{N}$), so $|AB| = \mathcal{O}(N \cdot N / \sqrt{N}) = \mathcal{O}(N^{3/2})$. This intermediate inflates the second shuffle substantially (see binary shuffle column in Table~\ref{tab:triple_relation}).

\subsection*{Spark SQL Join (\texttt{sql\_join.py})}

The SQL strategy registers $A$, $B$, and $C$ as temporary views and delegates the join entirely to Catalyst:

\begin{lstlisting}[caption=Spark SQL Triple Join]
result = spark.sql("""
    SELECT A.x, A.y, B.z, C.w
    FROM A JOIN B ON A.y = B.y
    JOIN C ON B.z = C.z
""")
\end{lstlisting}

Catalyst analyses relation statistics, selects the join order and algorithm (broadcast hash join, shuffle hash join, or sort-merge join), and inserts exchange operators automatically. With our synthetic equal-sized relations, Catalyst consistently selects a \texttt{BroadcastHashJoin} for the build-side relation, drastically reducing the shuffle footprint. The equi-join predicates remain inside the JVM throughout, eliminating the Python serialisation overhead that burdens the RDD strategies.

\section*{Experimental Setup}

Metrics are collected as $\Delta$ between Spark's cumulative accumulators before and after each strategy:
\begin{itemize}
    \item \textbf{Wall-clock time}: measured around the \texttt{.count()} action.
    \item \textbf{Task count}: total Spark tasks spawned across all stages.
    \item \textbf{Shuffle volume}: shuffle bytes written $+$ shuffle bytes read.
\end{itemize}
A 1-second sleep and cache clear between strategies prevents warm-cache interference.

Two experiments are conducted:
\begin{enumerate}
    \item \textbf{Varying relation size} (fixed $b = c = 24$): $N \in \{10{,}000;\; 50{,}000;\; 100{,}000;\; 200{,}000\}$.
    \item \textbf{Varying reducer count} (fixed $N = 100{,}000$): $b = c \in \{4,\; 12,\; 24,\; 36,\; 100\}$.
\end{enumerate}

\section*{Evaluation and Results}
\begin{table}[htbp]
\centering
\caption{Triple Join — Varying Relation Size ($b = c = 24$ reducer buckets per dimension)}
\label{tab:triple_relation}
\begin{tabular}{llrrrr}
\toprule
\textbf{$N$} & \textbf{Strategy} & \textbf{Rows} & \textbf{Time (s)} & \textbf{Tasks} & \textbf{Shuffle (B)} \\
\midrule
\multirow{3}{*}{10,000}
 & Ternary (Hypercube) & 99,000     & 2.370  & 96  & 1,706,316   \\
 & Binary              & 99,000     & 1.238  & 120 & 2,636,942   \\
 & Spark SQL           & 99,000     & 0.358  & 49  & 2,706       \\
\midrule
\multirow{3}{*}{50,000}
 & Ternary (Hypercube) & 2,475,000  & 4.495  & 96  & 7,248,954   \\
 & Binary              & 2,475,000  & 3.132  & 120 & 33,990,786  \\
 & Spark SQL           & 2,475,000  & 0.392  & 49  & 2,694       \\
\midrule
\multirow{3}{*}{100,000}
 & Ternary (Hypercube) & 9,900,000  & 9.726  & 96  & 53,141,488  \\
 & Binary              & 9,900,000  & 9.844  & 120 & 126,096,688 \\
 & Spark SQL           & 9,900,000  & 0.417  & 49  & 2,694       \\
\midrule
\multirow{3}{*}{200,000}
 & Ternary (Hypercube) & 39,600,000 & 46.464 & 96  & 29,582,018  \\
 & Binary              & 39,600,000 & 43.077 & 120 & 486,576,818 \\
 & Spark SQL           & 39,600,000 & 0.634  & 49  & 2,694       \\
\bottomrule
\end{tabular}
\end{table}


\begin{table}[htbp]
\centering
\caption{\textbf{Triple Join}: Varying Reducer Count ($N = 100{,}000$; $b = c$)}
\label{tab:triple_reducers}
\begin{tabular}{llrrrr}
\toprule
\textbf{Reducers ($b \times c$)} & \textbf{Strategy} & \textbf{Rows} & \textbf{Time (s)} & \textbf{Tasks} & \textbf{Shuffle (B)} \\
\midrule
\multirow{3}{*}{$4\times4=16$}
 & Ternary             & 9,900,000 & 11.650 & 16  & 11,554,986  \\
 & Binary              & 9,900,000 & 13.647 & 20  & 122,348,540 \\
 & Spark SQL           & 9,900,000 & 0.442  & 9   & 454         \\
\midrule
\multirow{3}{*}{$12\times12=144$}
 & Ternary             & 9,900,000 & 9.900  & 48  & 26,586,996  \\
 & Binary              & 9,900,000 & 10.016 & 60  & 125,074,972 \\
 & Spark SQL           & 9,900,000 & 0.411  & 25  & 1,350       \\
\midrule
\multirow{3}{*}{$24\times24=576$}
 & Ternary             & 9,900,000 & 9.696  & 96  & 53,141,488  \\
 & Binary              & 9,900,000 & 9.779  & 120 & 126,096,688 \\
 & Spark SQL           & 9,900,000 & 0.432  & 49  & 2,694       \\
\midrule
\multirow{3}{*}{$36\times36=1{,}296$}
 & Ternary             & 9,900,000 & 9.985  & 144 & 81,871,606  \\
 & Binary              & 9,900,000 & 9.781  & 180 & 127,483,636 \\
 & Spark SQL           & 9,900,000 & 0.480  & 73  & 4,038       \\
\midrule
\multirow{3}{*}{$100\times100=10{,}000$}
 & Ternary             & 9,900,000 & 18.100 & 400 & 255,793,296 \\
 & Binary              & 9,900,000 & 15.085 & 500 & 136,926,434 \\
 & Spark SQL           & 9,900,000 & 0.480  & 201 & 11,206      \\
\bottomrule
\end{tabular}
\end{table}

\paragraph{Correctness.}
Every run produces identical row counts (e.g.\ 9,900,000 at $N = 100{,}000$).

\paragraph{Spark SQL dominates by $\mathbf{20\times}$ or more.}
Across every configuration, SQL finishes in under 0.64\,s while both RDD strategies require 2–46\,s. The shuffle footprint of SQL is at most a few kilobytes — 2,694\,B even at $N = 100{,}000$, compared to tens or hundreds of megabytes for the RDD approaches. This is the same phenomenon observed in Part~1: the DataFrame/SQL API stays inside the JVM, and Catalyst rewrites the three-way join into a \texttt{BroadcastHashJoin} where the smaller build-side is broadcast to each executor without a full sort-merge shuffle. The near-constant 2,694\,B shuffle across $N \in \{50{,}000;\; 100{,}000\}$ indicates Catalyst successfully broadcasts at least one relation in every case, bypassing the network entirely for that side.

\paragraph{Ternary vs.\ Binary:}
At small-to-medium $N$ (up to $100{,}000$ with $b = c = 24$), the Ternary is marginally faster than Binary (9.70\,s vs.\ 9.78\,s) and produces less shuffle volume (53\,MB vs.\ 126\,MB), a factor-of-$2.4\times$ saving that aligns with the theoretical expectation. The single shuffle of the Hypercube moves $\sim N(2b+1)$ tagged records; Binary moves $\sim 2N$ records in the first shuffle \emph{plus} the intermediate $|AB|$ in the second. Because the join on $y$ has high fan-out in our synthetic data, $|AB| = \mathcal{O}(N^{3/2}/\sqrt{N}) = \mathcal{O}(N)$ rows with wide schemas, pushing the binary shuffle to 126\,MB vs.\ 53\,MB.

At $N = 200{,}000$, both RDD strategies degrade significantly (46\,s and 43\,s), but the ternary approach becomes \emph{slower} than binary despite producing less shuffle. At this scale the \texttt{groupByKey} boundary of the Hypercube must materialise all three tagged relations simultaneously in each reducer's memory — including $c$ replicas of every $A$-row and $b$ replicas of every $C$-row — causing GC pressure and spill-to-disk that offsets the reduced network cost. The binary strategy avoids this by externalising the intermediate to disk between the two cogroup shuffles.

\paragraph{Effect of reducer count on the Ternary.}
Table~\ref{tab:triple_reducers} reveals a clear sweet-spot at $b = c \approx 24$. With only $4 \times 4 = 16$ reducers, each partition handles $1/16$ of the data, but must materialise $c = 4$ replicas of each $A$-row and $b = 4$ replicas of each $C$-row per local join, producing large in-memory groups and high per-reducer CPU latency (11.65\,s). Increasing to $12 \times 12$ or $24 \times 24$ reducers improves throughput to $\approx 9.7$\,s by spreading the workload more finely. Pushing to $100 \times 100 = 10{,}000$ reducers reverses the trend: task-scheduling overhead and the explosion of shuffle files (400 reduce tasks, 256\,MB shuffled) drives ternary time back up to 18.1\,s. The shuffle \emph{volume} for Ternary grows monotonically with $b = c$ (from 11.6\,MB at $b=4$ to 255.8\,MB at $b=100$) because total replication cost $\propto |A| \cdot c + |C| \cdot b = 2bN$: more reducers mean more replicas, not fewer.

\paragraph{Effect of reducer count on Binary Joins.}
The Binary strategy's shuffle volume stays relatively stable across reducer counts (122--137\,MB for $b \in \{4,\ldots,36\}$) because the dominant cost is the intermediate $AB$ materialisation in the second shuffle, which is insensitive to the partition count. At $b = 100$, the binary shuffle rises only slightly to 137\,MB while ternary explodes to 256\,MB, causing the \emph{only} crossover: Binary outperforms Ternary (15.1\,s vs.\ 18.1\,s). This illustrates that the Hypercube's replication advantage holds only when $b$ and $c$ are calibrated to the available per-reducer memory and the expected join cardinality.

\end{document}
